\section{How to read the formal development}\label{sec:reading}
The development uses two complementary mathematical languages. Arithmetic packets begin with sequences and formal power series, then transport divisibility through exact coordinate changes. Polynomial lifting begins with a finite residual equation, passes to compatible digit towers, and finally enters complete p-adic and analytic spaces. The transition between these languages is explicit: a rational coefficient has to be shown integral before it becomes a p-adic integer; an integer congruence tower has to be shown Cauchy before it becomes a root; a formal polynomial derivative has to be identified with a strict Fr\'echet derivative before it supports an analytic inverse theorem.

This guide follows those transitions. Read the \emph{contract} before the excerpt: it states the coefficient domain, ambient dimensions, and hypotheses that may be declared earlier in the source. The excerpt is an exact contiguous selection from a preserved Lean file. Its line range and source hash are supplied in the accompanying excerpt inventory. The proof reading then explains what the displayed tactics accomplish mathematically. A theorem imported from another module is a dependency, not a hidden new assumption; a hypothesis supplied to a theorem remains a mathematical assumption even when its name does not appear in the theorem's axiom closure.

\subsection{Statements, constructions, audits, and identifiers}
The current consolidated manuscript links 62 statements to 202 named source declarations. That index resolves names and source identities. It does not prove that every informal sentence equals a Lean statement. Several declarations may establish one displayed theorem, while one source theorem may support several expository consequences. The restored Bernoulli specialization has a classical formal-series derivation in the manuscript, but no dedicated Lean bridge was matched. The manuscript's open problems retain their open status.

The full declaration index has 1,876 rows. A row records a declaration kind, qualified name, and entire source-file hash. Byte-identical copies are collapsed; different source versions remain separate. There are 1,863 distinct kind/name pairs, including 1,542 distinct names whose command kind is \code{theorem}. These are identifier counts, not counts of independent mathematical results. The rows also contain constructions, statement contracts, examples and control support. Thirteen names have more than one source version.

At the latest polynomial regularity checkpoint, the 55 proof modules contain 322 explicit theorem commands, 76 definitions and seven abbreviations. Eight positive control modules add 61 theorem commands and 14 definitions. Thirteen statement-contract modules add 32 theorem commands, two definitions and one abbreviation. The resulting 515 named declarations have matching retained axiom records. The frequently quoted number 415 is the total theorem-command count across these three groups. It includes the 61 control and 32 contract commands.

\subsection{Reading the Lean notation}
\code{Fin d → R} means a vector with exactly $d$ coordinates over $R$. In rectangular lifting, $d$ counts unknowns and $e$ counts equations. In a parameter split, \code{Fin (e+k) → R} separates $e$ solved coordinates and $k$ free parameters. A theorem using \code{Fin (m+k)} may instead have $m$ solved coordinates and $e$ equations; the distinction matters when a fiber has $p^{m-e}$ elements.

\code{A →ₗ[R] B} is a linear map, while \code{A →L[R] B} is a continuous linear map. An arrow decorated \code{≃L[R]} is a continuous linear equivalence. The analytic derivative is a continuous linear map; a raw matrix becomes one through the explicitly defined \code{matrixCLM}. \code{IsUnit J} for a square matrix means it is invertible in the matrix ring. It does not mean each matrix entry is a scalar unit.

\code{ZMod p} is the residue ring. The source uses \code{Fact p.Prime} when it needs a finite field, rank, or p-adic objects. Some affine and integer polynomial statements require only $p>0$ or $q>0$; their quotient algebra does not require primality. A hypothesis \code{1 ≤ L} supplies the depth needed for polynomial first-order linearization. The affine depth-zero case is valid because an affine map has no quadratic Taylor remainder.

\code{∃! b, P b} means there is exactly one object satisfying the full predicate $P$. A local root theorem's predicate includes its residue ball. It does not imply that the polynomial has only one root in the entire field. Likewise, an analytic graph theorem supplies open parameter and solved-coordinate neighborhoods; global uniqueness is outside its conclusion.

\subsection{What the retained proof evidence establishes}
The original compiler runs use Lean 4.34.0-rc2 and the recorded Mathlib revision \code{aeb72f1e9981b6d35b426b16dca7759be3e7ba06}. Their permitted axiom closure is \code{propext}, \code{Classical.choice}, and \code{Quot.sound}. This is a check for standard foundational dependencies and absence of a recorded admitted-proof axiom. It is not a claim that the proofs avoid classical choice, nor that every manuscript hypothesis has been discharged.

Fresh reconciliation reproduced all 15 retained evidence frontiers: 644 source-hash comparisons and 713 recorded log-hash comparisons match. Four older standalone affine logs have no digest in their original report; they are fingerprinted in the later index, so their historical provenance remains weaker. The fresh strict parser matched all 515 latest named axiom records to source declarations, and the unchanged Python parser regression rejected duplicated, mismatched, and forbidden-axiom records. These are static checks. Kernel checking remains the inherited compiler evidence.
